% Part: sets-functions-relations
% Chapter: size-of-sets-complete

\documentclass[../../../include/open-logic-chapter]{subfiles}

\begin{document}

\olchapter{sfr}{siz}{समुच्चयों का आकार}

\begin{editorial}
यह अध्याय गणनों, गणनीयता और अगणनीयता पर चर्चा करता है। कई खंडों के दो रूप
हैं: एक अधिक प्रारंभिक, जो गणनों को सूचियाँ या $\PosInt$ से आच्छादक फलन मानता
है; और दूसरा अधिक अमूर्त, जो गणनों को $\Nat$ के साथ एकैकी आच्छादक
प्रतिचित्रणों के रूप में परिभाषित करता है।
\end{editorial}

\olimport{introduction}

\olimport{enumerability}

\olimport{zig-zag}

\olimport{pairing}

\olimport{pairing-alt}

\olimport{non-enumerability}

\olimport{reduction}

\olimport{equinumerous-sets}

\olimport{comparing-size}

\olimport{schroder-bernstein}

\begin{editorial}
निम्न \olref[sfr][siz][enm-alt]{sec},
\olref[sfr][siz][nen-alt]{sec}, \olref[sfr][siz][red-alt]{sec},
\olref[sfr][siz][enm]{sec}, \olref[sfr][siz][nen]{sec} और
\olref[sfr][siz][red]{sec} के वैकल्पिक रूप हैं; इन्हें टिम बटन ने अपने ओपन
सेट थ्योरी पाठ के लिए लिखा है। वे कुछ अधिक उन्नत हैं और गणनीयता की ऐसी भिन्न
परिभाषा का उपयोग करते हैं जो समुच्चय-सिद्धांत के संदर्भ में अधिक उपयुक्त है
(अर्थात्, $\Nat$ या उसके किसी आरंभिक खंड के साथ एकैकी आच्छादक प्रतिचित्रण
होना, न कि सूचीबद्ध किए जा सकना या $\PosInt$ से किसी आच्छादक फलन का परिसर
होना)।
\end{editorial}

\olimport{enumerability-alt}
\olimport{non-enumerability-alt}
\olimport{reduction-alt}

\OLEndChapterHook

\end{document}
